\documentclass[11pt]{article}
\usepackage[a4paper,margin=27mm]{geometry}
\usepackage{amsmath,amssymb,amsthm,mathtools,mathrsfs}
\usepackage{graphicx}
\usepackage{microtype}
\usepackage[hidelinks]{hyperref}

\newtheorem{theorem}{Theorem}
\newcommand{\A}{\mathcal A}
\newcommand{\Hinf}{\mathcal H^\infty}
\newcommand{\B}{B_{\ell_2}}
\newcommand{\M}{\mathcal M}
\newcommand{\F}{\mathscr F}
\newcommand{\C}{\mathbb C}
\newcommand{\T}{\mathbb T}

\title{Coordinate-power middle fibers contain a full analytic ball\\
\large A solution of the highlighted example in arXiv:2102.06771}
\author{Candidate solution packet}
\date{August 11, 2026}

\begin{document}
\maketitle

\noindent\textbf{Status.}
Candidate full resolution of the highlighted coordinate-square example,
likely valid pending expert review.

\section{Source question}

Let $\B$ be the open unit ball of the complex Hilbert space $\ell_2$.
Dimant and Singer study the vector-valued spectrum
\[
 \M_{u,\infty}(\B,\B)
 =\{\Phi:\A_u(\B)\longrightarrow\Hinf(\B):
       \Phi\text{ is a nonzero algebra homomorphism}\}.
\]
It has the natural projection
\[
 \xi(\Phi)(x)=
 \bigl(\Phi(\langle\,\cdot,e_n\rangle)(x)\bigr)_{n\ge1}
 \quad\text{and fibers}\quad
 \F(g)=\xi^{-1}(g).
\]
Their Section~2 shows that some middle fibers are singletons while others
contain an analytic copy of a ball.  At the end of the section they ask what
happens for the natural middle symbol
\[
 g_2(x)=(x_n^2)_{n\ge1},
\]
which satisfies neither of their sufficient criteria \cite{DimantSinger}.

\begin{figure}[ht]
  \centering
  \includegraphics[width=\textwidth]{figures/open_problem_crop.png}
  \caption{The highlighted unresolved example at the end of Section 2,
  source PDF page 9.}
\end{figure}

The source passage asks for a description of this fiber and places it in the
broader question whether singleton fibers and fibers containing analytic balls
are the only possible middle-fiber behaviors.

\section{Proof intuition}

The obstruction to applying the source's ball criterion is that $g_2$ loses
no input direction.  Instead of looking for a kernel direction, one can find
norm room intrinsic to the coordinate-power map.  The fixed holomorphic
factor $\phi_m(x)=x_1\cdots x_m$ satisfies
\[
 \|g_m(x)\|^2+|\phi_m(x)|^2\le \|x\|^{2m},
 \qquad g_m(x)=(x_n^m)_n.
\]
Thus an arbitrary bounded holomorphic vector map can be multiplied by
$\phi_m$ and placed in coordinates escaping to infinity without leaving the
Hilbert ball.  Weak-star ultrafilter limits of the corresponding composition
homomorphisms have projection $g_m$.

The escaping coordinates might appear to erase the parameter.  They do not:
diagonal quadratic and cubic polynomials convert those tails into generating
functions for the coordinate squares and cubes of the parameter.  Their
boundary values determine all coefficients.  Squares and cubes together
remove the sign ambiguity and recover the parameter exactly.

\section{Main result}

For $k\ge1$, let $P_k$ be the orthogonal projection onto
$\operatorname{span}\{e_1,\dots,e_k\}$ and let
$S_k e_j=e_{k+j}$.

\begin{theorem}\label{thm:main}
For every integer $m\ge2$, put
\[
 g_m(x)=(x_n^m)_{n\ge1}\qquad(x\in\B).
\]
Then $\|g_m\|=1$, $g_m(\B)\subset\B$, and there is an analytic injection
\[
 \Psi_m:B_{\Hinf(\B,\ell_2)}\longrightarrow \F(g_m).
\]
Moreover $\Psi_m(0)=C_{g_m}$.  In particular, the source's highlighted fiber
over $(x_n^2)_n$ contains an analytic copy of the full open unit ball of
$\Hinf(\B,\ell_2)$.
\end{theorem}

\begin{proof}
Fix $m\ge2$ and define
\[
 \phi_m(x)=x_1x_2\cdots x_m.
\]
If $a_n=|x_n|^2$, the multinomial expansion (first for finite partial sums,
then by monotone convergence) gives
\begin{equation}\label{eq:room}
 \sum_{n=1}^\infty a_n^m+a_1a_2\cdots a_m
 \le \left(\sum_{n=1}^\infty a_n\right)^m.
\end{equation}
Indeed, the terms $a_n^m$ occur with coefficient one and the mixed monomial
$a_1\cdots a_m$ occurs with coefficient $m!$.  Consequently
\begin{equation}\label{eq:defect}
 \|g_m(x)\|^2+|\phi_m(x)|^2\le \|x\|^{2m}<1
 \qquad(x\in\B).
\end{equation}
This also shows that $g_m(\B)\subset\B$; testing points $re_1$ as $r\uparrow1$
shows that $\|g_m\|=1$.

For $h\in B_{\Hinf(\B,\ell_2)}$ and $k\ge1$, define
\[
 s_{h,k}(x)=P_k g_m(x)+\phi_m(x)S_kh(x).
\]
The two summands have orthogonal coordinate supports.  Hence, by
\eqref{eq:defect},
\[
 \|s_{h,k}(x)\|^2
 =\|P_kg_m(x)\|^2+|\phi_m(x)|^2\|h(x)\|^2
 \le \|g_m(x)\|^2+|\phi_m(x)|^2<1.
\]
Thus $s_{h,k}:\B\to\B$ is holomorphic and defines the composition
homomorphism
\[
 \Psi_{m,k}(h)=C_{s_{h,k}}\in\M_{u,\infty}(\B,\B).
\]

Fix a free ultrafilter $\mathscr U$ on $\mathbb N$ containing every cofinite
set, and set
\[
 \Psi_m(h)=\lim_{k\to\mathscr U}\Psi_{m,k}(h)
\]
in the weak-star compact vector-valued spectrum.  The limit is an algebra
homomorphism.  The standard bounded-holomorphic ultrafilter argument used in
\cite[Lemma~2.1, Theorem~2.3, Example~2.12]{DimantSinger} shows that
$\Psi_m$ is analytic: for each $f\in\A_u(\B)$ and $x\in\B$, the scalar maps
$h\mapsto f(s_{h,k}(x))$ are bounded holomorphic functions, and their
weak-star ultrafilter limit is holomorphic.

For every fixed $n$ and every $k\ge n$,
\[
 \Psi_{m,k}(h)(\langle\,\cdot,e_n\rangle)(x)=x_n^m.
\]
Therefore $\xi(\Psi_m(h))=g_m$, so the image of $\Psi_m$ lies in
$\F(g_m)$.

It remains to prove injectivity.  For $q\in\{2,3\}$ and $\lambda\in\T$,
consider the continuous homogeneous polynomial
\[
 F_{q,\lambda}(y)=\sum_{j=1}^\infty \lambda^j y_j^q
 \qquad(y\in\ell_2).
\]
It belongs to $\A_u(\B)$.  Put
$\chi(\lambda)=\lim_{k\to\mathscr U}\lambda^k$; since $|\lambda|=1$,
$|\chi(\lambda)|=1$.  Direct substitution and passage to the ultrafilter
limit give, for $x\in\B$,
\begin{align}\label{eq:separate}
 \Psi_m(h)(F_{q,\lambda})(x)
 &=\sum_{n=1}^\infty \lambda^n x_n^{mq}
   +\chi(\lambda)\phi_m(x)^q
      \sum_{j=1}^\infty\lambda^j h_j(x)^q.
\end{align}

Suppose $\Psi_m(h)=\Psi_m(h')$.  On the nonempty open set
\[
 V_m=\{x\in\B:\phi_m(x)\ne0\},
\]
equation~\eqref{eq:separate}, first with $q=2$ and then with $q=3$, implies
for every $\lambda\in\T$ that
\[
 \sum_{j\ge1}\lambda^j h_j(x)^q
 =\sum_{j\ge1}\lambda^j h'_j(x)^q.
\]
These power series have absolutely summable coefficients and are continuous
on the closed disk.  Uniqueness of their Taylor coefficients yields
\[
 h_j(x)^2=h'_j(x)^2,
 \qquad h_j(x)^3=h'_j(x)^3
 \quad(j\ge1).
\]
For complex numbers, equality of both the square and cube forces equality;
hence $h_j(x)=h'_j(x)$ for every $j$ and every $x\in V_m$.  Each coordinate
of $h-h'$ is holomorphic, so the identity theorem on the connected ball
$\B$ gives $h=h'$ everywhere.  Thus $\Psi_m$ is injective.

Finally, when $h=0$ one has $s_{0,k}=P_kg_m$.  For every $x\in\B$,
$P_kg_m(x)\to g_m(x)$ in norm.  Since every $f\in\A_u(\B)$ is uniformly
continuous, $f(P_kg_m(x))\to f(g_m(x))$; therefore
$\Psi_m(0)=C_{g_m}$.
\end{proof}

\section{Consequences and limitations}

Taking $m=2$ fully determines which of the two known large-scale behaviors
occurs for the source paper's singled-out example: its fiber is not a
singleton, but contains the same full analytic parameter ball appearing in
the source's large-fiber results.  The theorem simultaneously treats every
coordinate power $m\ge2$.

The broader classification question remains open.  This argument does not
show that every middle fiber is either a singleton or contains an analytic
ball, and therefore does not exclude a third behavior for other symbols.

\section{Verification and novelty status}

The proof was audited at the following points: the multinomial norm estimate;
strict inclusion $s_{h,k}(\B)\subset\B$; coordinatewise fiber projection;
membership of the quadratic and cubic tests in $\A_u(\B)$; nonvanishing of
the ultrafilter phase on $\T$; coefficient recovery; and extension from
$V_m$ by the identity theorem.  The only imported mechanism is the source
paper's own weak-star ultrafilter construction of analytic maps into the
spectrum.

On August 11, 2026, the local run indexes were searched by arXiv id, exact
title, the displayed coordinate-square example, and middle-fiber terminology.
Targeted web searches found the source arXiv record and its 2022 journal
publication, but no later answer to this example.  The search was bounded, so
novelty confidence is moderate pending a specialist citation review.

\begin{thebibliography}{9}

\bibitem{DimantSinger}
V. Dimant and J. Singer,
\emph{A look into homomorphisms between uniform algebras over a Hilbert
space},
Studia Math. \textbf{265} (2022), no.~1, 57--75;
arXiv:2102.06771; DOI: 10.4064/sm210219-10-6.

\end{thebibliography}

\end{document}
